% Short note on the jet algorithms appearing in the FCCJetBenchmarks figures:
% what each one actually computes, and which FastJet/FCCAnalyses call produces it.
%
% Build:  pdflatex jet_algorithms.tex
\documentclass[11pt,a4paper]{article}
\usepackage[margin=2.4cm]{geometry}
\usepackage{amsmath}
\usepackage{amssymb}   % \gtrsim
\usepackage{booktabs}
\usepackage{hyperref}
\usepackage[T1]{fontenc}

\newcommand{\ee}{$e^{+}e^{-}$}
\newcommand{\dij}{d_{ij}}
\newcommand{\diB}{d_{iB}}

\title{Jet algorithms used in the FCC jet benchmarks}
\author{FCCJetBenchmarks}
\date{\today}

\begin{document}
\maketitle

\noindent
All jets are built with FastJet through the \texttt{FCCAnalyses}
\texttt{JetClustering::clustering\_*} functors, from
\texttt{src/histmaker\_tools/jets.py}. Every algorithm below is a sequential
recombination algorithm: repeatedly find the smallest of all $\dij$ and $\diB$,
merge the pair $ij$ (or, if a $\diB$ is smallest, declare $i$ a jet / recombine
it with the beam), and repeat. They differ only in the distance measure and in
the termination rule.

Unless stated otherwise the recombination scheme is the $E$-scheme (four-momenta
are simply added, \texttt{recombination}${}=0$), and the input is the collection
of visible particles with $|\eta| < 2.56$: stable generator particles with
neutrinos removed for gen jets, particle-flow objects for reco jets.

\section*{The \ee{} distance measures}

The generalized \ee{} $k_T$ family (FastJet \texttt{ee\_genkt\_algorithm}) uses
angular, not rapidity, distances, which is the natural choice at a lepton
collider where there is no beam remnant to be boosted along:
\begin{align}
  \dij &= \min\!\left(E_i^{2p},\, E_j^{2p}\right)\,
          \frac{1 - \cos\theta_{ij}}{1 - \cos R}\,,
  &
  \diB &= E_i^{2p}\,.
  \label{eq:eegenkt}
\end{align}
The exponent $p$ selects the algorithm:
\[
  p = -1 \ \ (\text{anti-}k_T), \qquad
  p = 0 \ \ (\text{Cambridge/Aachen}), \qquad
  p = +1 \ \ (k_T).
\]
All three are run in these benchmarks, inclusively, over the same radius grid.
Two properties of Eq.~\eqref{eq:eegenkt} matter for reading the figures:

\begin{itemize}
  \item The factor $1/(1-\cos R)$ is a \emph{global constant}: it never reorders
        the $\dij$ among themselves. The entire $R$ dependence therefore enters
        through the comparison of $\dij$ against $\diB$, i.e.\ through which
        particles are split off into their own jet rather than merged. Small $R$
        means more, narrower jets.
  \item As $R \to \pi$ the beam term can no longer win, and \ee{} generalized
        $k_T$ with $p=+1$ degenerates into the Durham algorithm below. Measured
        on simulated 6-jet events, exclusive-to-$N$ clustering with $p=+1$ is
        already \emph{bit-identical} to Durham for $R \gtrsim 0.6$. This is why
        the $k_T$ family in these benchmarks is run \emph{inclusively}, where $R$
        genuinely changes the jets.
\end{itemize}

\section*{Algorithms appearing in the figures}

\paragraph{Durham (\ee{} $k_T$, exclusive).}
FastJet \texttt{ee\_kt\_algorithm}; the classic LEP jet definition. It is
Eq.~\eqref{eq:eegenkt} with $p=+1$ and no $R$ at all, i.e.\ a pure
\[
  y_{ij} = \frac{2\,\min(E_i^2, E_j^2)\,(1-\cos\theta_{ij})}{E_{\text{vis}}^2}\,,
\]
run in \emph{exclusive} mode: the clustering is stopped when exactly
$N$ jets remain, with $N$ fixed per process to the expected number of
final-state jets ($2$, $4$ or $6$; see \texttt{src/process\_config.py}).
Because $N$ is imposed, there is no jet-multiplicity distribution to speak of
and no radius to scan. Call:
\begin{center}
\texttt{JetClustering::clustering\_ee\_kt(2, N, 1, 0)}
\end{center}
(\texttt{exclusive}${}=2$ meaning ``exactly \texttt{cut} jets'', \texttt{cut}${}=N$,
$E$-ordered output).

\paragraph{\ee{} Cambridge/Aachen, inclusive, radius scan.}
Eq.~\eqref{eq:eegenkt} with $p=0$, so $\dij$ is purely angular:
$\dij = (1-\cos\theta_{ij})/(1-\cos R)$, independent of energy. Pairs are
merged in order of increasing opening angle. Run inclusively, so the number of
jets is an output, not an input. Scanned over
$R = 0.4, 0.6, 0.8, 1.0, 1.2, 1.4$. Call:
\begin{center}
\texttt{JetClustering::clustering\_ee\_genkt(R, 0, 0, 0, 0, 0)}
\end{center}
Method directories \texttt{PF\_EECambridgeR04}\,\dots\,\texttt{PF\_EECambridgeR14}.

\paragraph{\ee{} $k_T$, inclusive, radius scan.}
Eq.~\eqref{eq:eegenkt} with $p=+1$: $\dij$ is weighted by the \emph{softer}
of the two energies, so soft particles are clustered first and jets grow from
the soft radiation inwards. Same inclusive mode and same radius grid as the C/A
scan, which makes the two directly comparable at matched $R$. Call:
\begin{center}
\texttt{JetClustering::clustering\_ee\_genkt(R, 0, 0, 0, 0, 1)}
\end{center}
Method directories \texttt{PF\_EEKtR04}\,\dots\,\texttt{PF\_EEKtR14}.
At equal $R$, $k_T$ merges considerably more aggressively than C/A: at $R=0.8$ in
$Z(\to qq)H(\to WW \to qqqq)$, $74\%$ of events still yield $6$ jets under $k_T$
against $98\%$ under C/A.

\paragraph{\ee{} anti-$k_T$, inclusive, radius scan.}
Eq.~\eqref{eq:eegenkt} with $p=-1$: $\dij$ is weighted by the \emph{inverse} of the
harder energy, so hard particles cluster first and soft radiation is swept into the
nearest hard core rather than forming jets of its own. Same inclusive mode and radius
grid as the other two, making all three comparable at matched $R$. Call:
\begin{center}
\texttt{JetClustering::clustering\_ee\_genkt(R, 0, 0, 0, 0, -1)}
\end{center}
Method directories \texttt{PF\_EEAntiKtR04}\,\dots\,\texttt{PF\_EEAntiKtR14}.
First produced in September 2026; before that the repository contained no genuine
anti-$k_T$ run at all (see the naming history below).

\paragraph{Ordering of the three exponents.}
Merging becomes monotonically more aggressive as $p$ increases. Measured at $R=0.8$ in
$Z(\to qq)H(\to WW \to qqqq)$ over $10^4$ events:

\begin{center}
\small
\begin{tabular}{@{}rlcc@{}}
\toprule
$p$ & algorithm & events yielding 6 jets & fully-matched pass rate \\
\midrule
$-1$ & anti-$k_T$ & $99.8\%$ & $0.878$ \\
$0$ & Cambridge/Aachen & $98.4\%$ & $0.815$ \\
$+1$ & $k_T$ & $74.4\%$ & $0.611$ \\
\midrule
--- & Durham (exclusive $N$) & $100\%$ by construction & $0.866$ \\
\bottomrule
\end{tabular}
\end{center}

\paragraph{Energy recovery (the ``-ER'' variants).}
Not a clustering algorithm but a post-processing step applied to the inclusive
C/A output: the $N$ leading jets are kept and every surplus jet is recombined
into them, so the event ends with exactly $N$ jets
(\texttt{ZHfunctions::energy\_recovery}, \texttt{src/histmaker\_functions/functions.h}).
It exists to recover the energy that inclusive clustering leaves in extra jets,
and is available for all three exponents
(\texttt{PF\_E\_recovery\_EECambridgeR*}, \texttt{PF\_E\_recovery\_EEKtR*},
\texttt{PF\_E\_recovery\_EEAntiKtR*}).

Because it acts only where the clustering overproduces, its size is governed by
$N$. Measured at $R = 0.4$ on $Z(\to\nu\nu)H(\to qq)$ ($N = 2$) it moves the
$m_H$ peak from $119.9$ to $124.9$~GeV, and makes the result independent of both
the radius and the exponent --- the $N$ jets then contain every visible particle,
so nothing is left for the jet definition to change. On the 6-jet light-flavour
process at $R \ge 0.8$, by contrast, the clustering already yields $\le N$ jets
and energy recovery is exactly the identity.

\paragraph{CaloJets.}
Not clustered in this framework at all: the jets are read from the
\texttt{CaloJetDurham} branch produced by Delphes, which applies the Durham
algorithm to calorimeter towers rather than to particle-flow objects. Included
as the ``no particle flow'' reference.

\paragraph{Hadron-collider anti-$k_T$.}
Available as \texttt{--jet-algorithm AK}, i.e.\ \texttt{clustering\_antikt}
(FastJet \texttt{antikt\_algorithm}), with the usual $p_T$/rapidity measure
\[
  \dij = \min\!\left(p_{T,i}^{-2},\,p_{T,j}^{-2}\right)
         \frac{\Delta R_{ij}^2}{R^2}\,,
  \qquad
  \diB = p_{T,i}^{-2}\,.
\]
It is not used in any of the production figures and is kept only for
cross-checks: the $p_T$/rapidity measure is not the appropriate one at a lepton
collider.

\section*{Summary}

\begin{center}
\small
\begin{tabular}{@{}llccl@{}}
\toprule
Label in figures & FastJet & $p$ & Mode & Method directory \\
\midrule
Durham / PFJets & \texttt{ee\_kt} & $+1$ & excl., $N$ jets & \texttt{PF\_Durham} \\
CaloJets & \texttt{ee\_kt} (Delphes) & $+1$ & excl., $N$ jets & \texttt{CaloJets\_Durham} \\
ee-C/A $R$ & \texttt{ee\_genkt} & $0$ & inclusive & \texttt{PF\_EECambridgeR*} \\
ee-C/A $R$ -ER & \texttt{ee\_genkt} & $0$ & incl.\ $+$ recovery & \texttt{PF\_E\_recovery\_EECambridgeR*} \\
ee-$k_T$ $R$ & \texttt{ee\_genkt} & $+1$ & inclusive & \texttt{PF\_EEKtR*} \\
ee-anti-$k_T$ $R$ & \texttt{ee\_genkt} & $-1$ & inclusive & \texttt{PF\_EEAntiKtR*} \\
\bottomrule
\end{tabular}
\end{center}

\section*{Naming history (important when reading older figures)}

Until September 2026 the C/A scan was labelled ``generalized \ee{} anti-$k_T$''
and its output directories were named \texttt{PF\_AntiKtR*}. That label was
wrong. The clustering call was
\begin{center}
\texttt{JetClustering::clustering\_ee\_genkt(R, 0, 0, 0)}
\end{center}
with only four arguments, whereas the constructor signature is
\begin{center}
\texttt{(radius, exclusive, cut, sorted, recombination=0, exponent=0.)}
\end{center}
The exponent therefore defaulted to $0$, so every one of those runs was
Cambridge/Aachen and never anti-$k_T$. The histograms themselves are valid C/A
data, so they were relabelled rather than re-run, and the directories renamed to
\texttt{PF\_EECambridgeR*}.

Genuine \ee{} anti-$k_T$ ($p=-1$) is now run separately, as \texttt{EEAKT} with output
in \texttt{PF\_EEAntiKtR*}. Note that \texttt{EEAK} (no trailing \texttt{T}) and
\texttt{EEAKT} are different algorithms: the former is the historical spelling that ran
$p=0$ and survives only as a deprecated alias of \texttt{EECA}. Likewise
\texttt{PF\_AntiKtR*} (Cambridge/Aachen, legacy) must not be confused with
\texttt{PF\_EEAntiKtR*} (genuine anti-$k_T$).

Any figure or table showing ``ee-AK'', ``anti-kt'' or \texttt{AntiKtR} should be
read as \ee{} Cambridge/Aachen.

\end{document}
